\clearpage
\hypertarget{native-469-letter}{}
\begin{center}
\rule{0.55\linewidth}{0.4pt}\\[2pt]
Astronomisehe Nachrichten. Band XXI. Nr. 495. S. 235..238. 8. Februar 1844.\\[2pt]
\rule{0.55\linewidth}{0.4pt}
\end{center}

\begin{center}[\textsc{Gauss} an \textsc{Schumacher.}]\quad Göttingen, 17. Januar 1844.\end{center}

Anbei sende ich Ihnen eine Ephemeride des \emph{Kometen}, die Herr Dr. \textsc{Goldschmidt} nach seinen zweiten Elementen berechnet hat. Ich zweifle nicht daran, dass man den Kometen damit immer leicht finden kann, sofern sein Licht nicht so schwach wird, dass er nicht mehr zu sehen ist.

\Needspace{3\baselineskip}
\SourceLeaf{470}
Herr Dr. \textsc{Goldschmidt} wird den Vergleich dieser zweiten Elemente mit sämtlichen ihm bekannt gewordenen Beobachtungen unverzüglich abschließen.

Meine eigene Beobachtung vom 11. Januar lautet:
\begin{center}
\begin{tabular}{l|l|l|l}
1844 Januar 11 & $9^{\rm u}23^{\rm m}41^{\rm s}.8$ & $77^\circ8'49''4$ & ---\quad ---\quad --- \\
 & $9\;23\;47.5$ & ---\quad ---\quad --- & $+3^\circ32'26''7$ \\
\end{tabular}
\end{center}

Die Deklination bleibt etwas ungewiss. Der geringe Unterschied zwischen den Zeiten, für die Rektaszension und Deklination gelten, rührt daher, dass die Rektaszension auf drei Beobachtungen beruht, die Deklination auf zwei, und dass die dritte Beobachtung, die nur für die Rektaszension brauchbar ist, \emph{nahe}, aber \emph{nicht genau}, in der Mitte zwischen den beiden anderen liegt.

Ich habe die Chronometer eines Handwerkers aus Hannover jetzt hier zur Prüfung, allerdings erst seit 4 oder 5 Tagen. In dieser Zeit war ihr Gang ausgezeichnet gleichmäßig, obwohl sich in so kurzer Zeit noch wenig entscheiden lässt.

Schließlich merke ich noch an, dass sich meine Beobachtung vom 13. Dezember bei der neuen Reduktion geringfügig geändert hat, allein wegen eines Fehlers bei der \emph{Reduktion des Sterns}. Sie lautet jetzt:
\begin{center}
\begin{tabular}{l|l|l|l}
1842 Dezember 13 & $9^{\rm u}35^{\rm m}15^{\rm s}.9$ & ---\quad ---\quad --- & $+3^\circ39'28''6$ \\
 & $9\;43\;13.7$ & $78^\circ45'38''0$ & ---\quad ---\quad --- \\
\end{tabular}
\end{center}

